\chapter{Theoretical Background and Related Work}\label{ch:background}

\section{Conceptual and Regulatory Background}

\subsection{Credit Scoring Models and Interpretability}

Credit scoring is the task of estimating the probability that a borrower will
default, framed here as a binary classification problem \citep{lessmannBenchmarkingStateoftheartClassification2015}. Models
are grouped in this thesis by a structural criterion: whether the fitted object can
be read directly, or whether a second computation is needed to say what it does. The
criterion concerns the availability of an explanation, not its quality. An
intrinsically interpretable model with hundreds of correlated or transformed features
can be hard to grasp.

Logistic regression satisfies the criterion structurally because its coefficients
determine the prediction directly
\citep{liptonMythosModelInterpretability2018}. Their interpretation nevertheless
depends on feature scaling, categorical coding and the correlation structure, and
they act on the log-odds rather than directly on the probability.

The EBM provides a non-linear yet component-wise readable \emph{glass-box}
alternative. It is a generalised additive model extended by a limited number of
pairwise interactions, a GA\textsuperscript{2}M in the lineage of
\citet{louAccurateIntelligibleModels2013} and
\citet{noriInterpretMLUnifiedFramework2019}. It breaks down the link-transformed
expectation of default into
\begin{equation}
g\!\left(\mathbb{E}[y \mid x]\right) \;=\; \beta_0 \;+\; \sum_{j} f_j(x_j)
\;+\!\! \sum_{(j,k)\,\in\,\mathcal{I}} \!\! f_{jk}(x_j, x_k),
\label{eq:ebm}
\end{equation}
where $g$ is the link function, the logit here, so that the terms act on the log-odds
as in logistic regression. Each $f_j$ is a piecewise constant step function over the
binned values of feature~$j$, and $\mathcal{I}$ is the set of selected
interaction terms. Every term can
be plotted and read off exactly, in the sense that the terms reproduce the model
output rather than approximate it. They describe learned associations, not causal
effects \citep{amekoeExploringAccuracyInterpretability2024b}.

Gradient-boosted tree ensembles are constructed the other way round, building shallow
trees in stages, each fitted to the \emph{negative} gradient of the loss with respect to
the current ensemble prediction, the pseudo-residuals \citep{friedmanGreedyFunctionApproximation2001}. XGBoost \citep{chenXGBoostScalableTree2016} adds penalties on
leaf weights and tree complexity; CatBoost \citep{prokhorenkovaCatBoostUnbiasedBoosting2018} grows symmetric
trees and offers ordered target statistics for categorical variables. A single such
tree would be readable, an
ensemble of several hundred is not, because the decision logic is distributed across
them. An attribution method must therefore be applied on top \citep{zouInterpretableCreditScoring2025}.

This thesis accordingly distinguishes intrinsic explanations, read directly from the
fitted logistic regression and EBM, from post-hoc attributions computed for XGBoost and
CatBoost.



\subsection{Profit-Based Evaluation}\label{sec:emp}

ROC-AUC evaluates how well a model ranks applicants but does not capture
the asymmetric economic consequences of lending decisions. Rejecting an
applicant who would repay forfeits the return on the loan, whereas granting
credit to an applicant who defaults causes a loss on the principal.
Cost-sensitive measures incorporate these consequences but require economic
parameters that are rarely known with certainty and may affect the resulting
ranking of models
\citep{verbrakenDevelopmentApplicationConsumer2014}.

The Expected Maximum Profit (EMP) measure accounts for this uncertainty by
averaging the maximum attainable return over a distribution of loss severities.
\citet{verbrakenNovelProfitMaximizing2013a} introduce the general framework,
which \citet{verbrakenDevelopmentApplicationConsumer2014} adapt to consumer
credit scoring. Let $\pi_D$ and $\pi_N$ denote the proportions of defaulters
and non-defaulters, and let $R_D(t)$ and $R_N(t)$ denote the respective
fractions rejected at cut-off $t$. The loss fraction associated with a default
is $\lambda$, the loss given default expressed as a fraction of the principal.
The return forgone by rejecting a non-defaulting applicant is represented by
$\mathrm{ROI}$. Profit per applicant and per unit of principal is therefore

\begin{equation}
P(t;\lambda,\mathrm{ROI})
= \lambda\,\pi_D R_D(t)
- \mathrm{ROI}\,\pi_N R_N(t).
\label{eq:profit}
\end{equation}

Profit is measured relative to granting all applications. Rejecting nobody
sets both rejection rates to zero; because this policy remains available,
the maximised profit cannot be negative.

Rather than fixing $\lambda$ at a single value, EMP integrates the maximum
attainable return over its mixed distribution $H$:

\begin{equation}
\mathrm{EMP}
= \int_{[0,1]}
P\bigl(t^*(\lambda);\lambda,\mathrm{ROI}\bigr)\,
\mathrm{d}H(\lambda),
\qquad
t^*(\lambda)
= \arg\max_t P(t;\lambda,\mathrm{ROI}).
\label{eq:emp}
\end{equation}

The loss-given-default distribution and the assumed return are specified in the
empirical design. EMP remains conditional on these assumptions and on the class
prevalence and should not be interpreted as realised portfolio profit. Absolute EMP
values are therefore compared between models within one dataset rather than across
datasets.

\subsection{Explainability in Machine Learning}\label{sec:xai}

Explainability methods are commonly organised along two axes: whether an
explanation is \emph{inherent} to the model or produced \emph{post hoc}, and
whether it is \emph{global} (describing the whole model) or \emph{local}
(describing a single decision) \citep{nautaAnecdotalEvidenceQuantitative2023}. \citet{liptonMythosModelInterpretability2018} further
decomposes interpretability into simulatability, decomposability and algorithmic
transparency, and stresses that a post-hoc explanation is not the same as
transparency. 

Additive feature contributions provide a common representation over the four
evaluated pipelines. \citet{lundbergUnifiedApproachInterpreting2017} derive SHapley Additive exPlanations (SHAP) from the Shapley value in cooperative game theory:
features are treated as players, the prediction as the payout, and each attribution
as a feature's average marginal contribution across feature coalitions. The resulting
explanation is additive: the attributions and a base value reconstruct the prediction.
Within the class of additive feature-attribution methods, Shapley values uniquely
satisfy local accuracy, missingness and consistency.

Several limitations qualify what such an attribution establishes. It describes the
\emph{model}, not the world: a large attribution indicates that the fitted model relies
on a feature, not that the feature causes default
\citep{rudinStopExplainingBlack2019a,
liptonMythosModelInterpretability2018}. Attributions are also defined relative to a
reference distribution, so the same prediction can receive different attributions
under different references
\citep{nautaAnecdotalEvidenceQuantitative2023}. With correlated features, the
allocation of attribution among them additionally depends on the assumed dependence
structure \citep{aasExplainingIndividualPredictions2021}.

Two quantitative properties recur in this literature: \emph{stability}, the
reproducibility of explanations under resampling or retraining, and
\emph{faithfulness}, the extent to which an explanation accurately reflects the
behaviour of the fitted model
\citep{nautaAnecdotalEvidenceQuantitative2023}.

Separately from these properties of individual attribution procedures,
\citet{rudinStopExplainingBlack2019a} argues that high-stakes settings should prefer
inherently interpretable models over black boxes explained post hoc. This is a
structural argument about the relation between the model and its explanation, but it does
not establish that every post-hoc attribution is numerically inaccurate or that every
intrinsically interpretable model is readily understandable.

\subsection{Regulatory Framework}\label{sec:regulatory}

Credit scoring in the European Union is governed principally by the AI Act and the
GDPR \citep{mullerAIbasedCreditScoring2025}. The AI Act
(Regulation~(EU)~2024/1689) first requires a system to fall within the definition of
an AI system under Article~3(1) \citep{EUAIAct2024}. Whether conventional statistical
scorecards satisfy that definition remains relevant: the Commission guidance
discussed by \citet{mullerAIbasedCreditScoring2025} treats well-established linear and
logistic regression as outside the definition, but the guidance is not legally
binding.

Assuming that a deployed scoring system falls within Article~3(1),
creditworthiness assessment is listed in Annex~III(5)(b). Because the use considered
here profiles natural persons and materially informs the credit decision, the
derogation in Article~6(3) would not apply. The classification therefore depends on
the system along with its intended use rather than on the selected model family.

For covered high-risk systems, the Act allocates obligations between providers and
deployers. Article~13 specifies transparency plus information requirements for the
system and its instructions for use, while Article~16(a) requires the provider to
ensure that these requirements are met. Article~15 requires an appropriate level of
accuracy and robustness. Article~86 provides a right to an explanation against the
deployer only under its statutory conditions: the decision must be based on the
system's output and adversely affect the person in a legally or similarly significant
way, and the right applies only insofar as it does not already follow from other Union
law. GDPR Article~15(1)(h), by contrast, concerns the controller \citep{EUAIAct2024, mullerAIbasedCreditScoring2025, GDPR2016}.

The GDPR applies independently. Article~22 provides a right, subject to exceptions,
not to be subject to a decision based solely on automated processing that produces
legal or similarly significant effects \citep{GDPR2016}. In
\emph{SCHUFA Holding} (C-634/21), the Court held that generating a credit score can
itself constitute such a decision where a third party draws strongly on it
\citep{OQGegenLand2023, mullerAIbasedCreditScoring2025}.

In \citet{CJEU2025DunBradstreet}, in the context of automated decision-making
within Article~22(1), the Court held that Article~15(1)(h) requires an intelligible
account of the procedure actually applied; disclosure of the algorithm or formula
alone is insufficient. Trade-secret
protection does not automatically preclude access. Where protected information is
involved, it must be made available to the competent authority or court for the
required balancing of interests.

The relevant high-risk requirements for Annex-III systems apply from
2~December~2027 under Regulation~(EU)~2026/1744
\citep{EUAIOmnibus2026}. The Act prescribes no numerical threshold for the
empirical metrics used in this thesis. ROC-AUC, calibration measures and attribution
analyses may constitute technical evidence, but they establish neither compliance nor
the legal sufficiency of an explanation.

Because this thesis evaluates models offline rather than assessing a deployed AI
system, the regulatory provisions serve as an interpretive benchmark. The analysis
does not constitute a conformity assessment.

\section{Related Work}

\subsection{Predictive and Economic Performance in Credit Scoring}

The benchmark study of
\citet{lessmannBenchmarkingStateoftheartClassification2015} found that
gradient-boosted and ensemble classifiers tend to outperform logistic regression in
discrimination across many credit datasets, although the advantage depends on the
dataset and evaluation metric. Reviewing 74 primary studies published between 2010
and 2018, \citet{dastileStatisticalMachineLearning2020a} similarly report that
boosting is the most frequently used and among the strongest-performing ensemble
approaches in credit scoring.

More recent work indicates that the practical difference can be smaller. Additive
models with interactions such as the EBM achieve predictive performance competitive
with black-box ensembles on tabular data
\citep{changHowInterpretableTrustworthy2021}. On German Credit, an additive boosting
model restricted to single-feature as well as pairwise terms slightly outperforms an
unconstrained XGBoost configuration while remaining inherently interpretable
\citep{zouInterpretableCreditScoring2025}. The best black-box and interpretable
models examined by \citet{dessainCostExplainabilityAI2023} differ by fifteen to twenty
basis points of return. Across 45 tabular datasets,
\citet{amekoeExploringAccuracyInterpretability2024b} report an average relative
accuracy loss below 4\% for intelligible models such as the EBM; this evidence
concerns tabular learning broadly rather than credit scoring specifically.

The work closest to this thesis is the preprint by
\citet{schwartzEnhancingMLModels2025}, which compares interpretable models with
XGBoost on Lending Club and reports similar observed performance for the EBM and
XGBoost in the settings examined. It uses SHAP only for feature selection and reports
neither explanation stability nor a profit metric.

A second strand evaluates credit-scoring models economically rather than
statistically. \citet{devosPROFITMAXIMIZINGLOGISTIC2018} fit logistic regression under
a profit-maximising objective,
\citet{chroscickiAdvantageCaseTailoredInformation2022} examine model development
against calculated profit rather than statistical criteria alone, and
\citet{xuProfitRiskdrivenCredit2024} optimise profit and risk jointly while treating
the cost parameters themselves as uncertain.

\subsection{Explainability Methods and Stability}

\citet{hlongwaneNovelFrameworkEnhancing2024} examine how SHAP can support the
interpretation of traditional scorecards. Moving beyond the presentation of feature
attributions alone, \citet{amekoeExploringAccuracyInterpretability2024b} include
\emph{stability}, the consistency of explanations under changes that should not
materially alter them, and \emph{faithfulness}, the extent to which an explanation
reflects the behaviour of the fitted model, among their evaluation criteria. Together,
these studies illustrate that explanation evaluation is multidimensional and cannot
be reduced to a single attribution output.

\citet{ballegeerEvaluatingStabilityModel2025b} provide a multi-factor analysis of
explanation stability in credit scoring, examining how cost-sensitive objectives,
class imbalance, learner and explainer choices, and model retraining affect SHAP and
LIME explanations. Their study provides evidence that explanation stability depends
on the wider modelling-and-explanation pipeline rather than on the explainer in
isolation. 

\subsection{Research Gap}\label{sec:gap}

The central research gap concerns the empirical status of the presumed
performance--explainability trade-off in credit scoring. \citet{bahloolPerformanceFairnessExplainability2026} find that predictive
performance, fairness and explainability are frequently studied in isolation and
that the trade-off between performance and explainability is more often assumed
than empirically quantified. Much of the literature consequently starts from a
comparatively complex predictive model and subsequently examines how its outputs
can be explained post hoc. Explainability is thereby commonly treated
as an additional layer following model selection rather than as a consideration
that may inform the initial choice of model \citep{bahloolPerformanceFairnessExplainability2026}.

This leaves an earlier model-choice question unresolved. The relevant issue is not
only how to explain a boosted-tree ensemble after it has been selected, but whether
a model whose fitted prediction function is directly inspectable can itself
constitute a viable scoring choice
\citep[see][]{rudinStopExplainingBlack2019a}. The
EBM makes this question empirically relevant because it can represent nonlinear
feature effects and selected pairwise interactions while retaining an additive
glass-box structure
\citep{louAccurateIntelligibleModels2013,
noriInterpretMLUnifiedFramework2019}. It therefore permits examination of the
presumed trade-off without equating inherent interpretability with a strictly
linear functional form.

Existing comparisons provide initial support for this model-choice perspective.
\citet{schwartzEnhancingMLModels2025} and
\citet{amekoeExploringAccuracyInterpretability2024b} report similar observed
predictive performance for the EBM and more opaque alternatives in their respective
settings, supporting its consideration as a viable model option rather than merely
a transparency-oriented substitute. Read in light of
\citet{bahloolPerformanceFairnessExplainability2026}, however, these outcomes do
not settle the presumed trade-off. Instead, they motivate closer examination of
whether similar observed performance recurs across datasets and extends to
profit-based performance. The remaining gap is therefore not the complete absence of
promising evidence for the EBM, but the limited basis for judging whether that
evidence continues across empirical settings and evaluation dimensions.

